\input{estilo}
\begin{document}
\begin{figura}\begin{tikzpicture}
  \anillo
  % los tres servidores de antes
  \foreach \ang/\n/\pos in {45/0/south west, -95/1/north west, 150/2/south east} {
    \nodoanillo{\ang}{\var{S}\textsubscript{\n}}{\pos}
  }
  % el nuevo, entre S1 y S2
  \fill[acento] (205:\R) circle (2.8pt);
  \node[etiq, acento, anchor=west] at (209:\R-0.15) {\var{S}\textsubscript{3}};
  % el arco viejo de S2, partido en dos
  \draw[->] (-100:\R+0.4) arc[start angle=-100, end angle=-148, radius=\R+0.4];
  \draw[->] (196:\R+0.4) arc[start angle=196, end angle=158, radius=\R+0.4];
  % la transferencia, desde un único vecino
  \draw[->, acento, dashed] (150:\R-0.35) arc[start angle=150, end angle=200, radius=\R-0.35]
    node[etiq, midway, right] {claves};
\end{tikzpicture}\end{figura}
\end{document}
